%=====================================================================
\section{Introduction}\label{sec:intro}

% Plan: the setting and why it is new: agents cross organisational boundaries through tools; every step has its own owner.
Software agents built on large language models increasingly act on behalf of people by calling tools, and the Model Context Protocol (MCP) has become a common way to expose those tools~\cite{mcp2025spec,hou2026mcp}. A single request may involve a person in one organisation, an agent platform operated by a second, a tool hosted by a third and a data platform owned by a fourth. Each of them governs its own step: the agent platform decides who may call which tool, the data platform decides which identity may read which resource, and the data owner decides for which purpose and towards which destination the data may be used. These decisions are written by different people, in different languages, and they change continually: tools are added and re-described, rules are edited, usage terms are tightened, and agents delegate to other agents.

% Plan: why it is hard; state the open issues documented in 2026.
Recent analyses agree that this composition is poorly governed. The authentication and authorisation of MCP lacks a standard framework across clients and servers~\cite{hou2026mcp}, authorisation is optional in the protocol and often absent in implementations~\cite{nsa2026mcp}, and current identity patterns serve one agent inside one trust domain rather than chains across domains~\cite{openid2025agentic}. Tool descriptions change after they were approved~\cite{bharti2026drift} and can carry instructions that steer agents~\cite{wang2026mcptox}. A recent survey of agent authorisation lists as open the mapping of one organisation's policy onto another's, re-approval when tools change, and attenuation along delegation chains~\cite{surapani2026authz}. Underneath these reports lies a problem that policy analysis has not addressed: \emph{whether the policies of several owners, taken together, still govern every call after each change, when no owner may see the others' policies.}

% Plan: why existing tools do not solve it.
Existing policy analysers answer a different question. They take one readable policy set and check it for conflicts, redundancies or properties~\cite{fisler2005margrave,hu2013anomalies,lin2010exam,backes2018zelkova,barnett2025authz,cutler2024cedar}. Work on multi-domain policies integrates or reconciles policies that are disclosed to an integrator~\cite{bonatti2002algebra,mazzoleni2008integration,mcdaniel2006reconciliation}. Privacy-preserving analyses exist for firewall rules and role constraints between two domains~\cite{chen2013firewall,liu2020smer} and for the evaluation of one request against several owners' policies~\cite{sheikhalishahi2019multiparty}, but none of them covers tools, usage terms, delegation or change. Re-running a global analysis after every change is also not an option: it is quadratic in the number of rules, it requires the very disclosure that owners refuse, and deciding even whether one layer covers every request is coNP-complete (Theorem~\ref{thm:hard}).

% Plan: our idea in one paragraph.
This paper takes a different view. Governance is an \emph{invariant} of the joint state, and the object to check is the \emph{change}. We model the rules of all owners and layers, the tool descriptions and the agent requests as hyperedges of one hypergraph over shared attribute-value vertices. The hypergraph is \emph{federated}, because each owner holds its own fragment, and \emph{dynamic}, because every modification is an event on it. A change can only affect the findings that involve the vertices it touches, so its consequences can be computed exactly in a small region, by the owners of that region only, on secret shares. We call the approach \name.

% Plan: contributions as bullets with forward references.
\paragraph{Contributions}
\begin{enumerate}
\item \textbf{Model.} A federated dynamic policy hypergraph that represents tool, data and usage rules of several owners, tool descriptions with their approved and current versions, agent requests and delegation chains in one structure (Section~\ref{sec:model}).
\item \textbf{Invariant and anomalies.} Six cross-owner anomalies, a duality between blocked grants and amplifications that holds for any combining function, attenuation along chains, and the hardness of universal coverage (Section~\ref{sec:model}).
\item \textbf{Admission of change.} A locality theorem that bounds the effect of every event to a vertex region, and an admission check proved necessary and sufficient for preserving the invariant, with a commutation result for concurrent events (Section~\ref{sec:admission}).
\item \textbf{Flags-only protocol.} A two-server secret-sharing protocol in which same-layer checks are inner products of incidence rows and cross-layer checks are products of bits computed by each owner on its own policy; a simulation-based security statement with an explicit leakage function, and an analysis of what the findings themselves reveal (Section~\ref{sec:protocol}).
\item \textbf{Ranking.} A federated hyperedge model that orders held changes by the expected use of the requests they affect (Section~\ref{sec:ranking}).
\item \textbf{Evaluation} on one year of real AWS managed-policy changes, 229,827 real MCP tool schemas, constructed multi-organisation federations, three benchmarks of an earlier sketch-based design, and real access requests (Section~\ref{sec:evaluation}).
\end{enumerate}

% Plan: main findings, honestly stated.
\paragraph{Main findings} On the AWS history, the region of a change covers a median of \awsRegMed\% of all statements, and in \awsEvents{} changes no finding appeared or disappeared outside it. Admission would have held \awsHeld{} changes (\awsHeldPct\%), among them \awsAmpROEv{} that gave agent-facing policies state-changing actions not held by the read-only job function, \awsAmpROAct{} actions in total. On constructed federations, incremental admission is exact on \eOneEvents{} random events, and every one of the \eOneSP{} pairwise and \eOneSC{} cross-layer secure checks reveals exactly the plaintext findings. The learned ranking helps only where partners know a resource that the organisation does not: AUROC \rkFedNew{} against 0.5 on such resources, and no gain on known ones. Section~\ref{sec:discussion} states what the approach does not do.

\paragraph{Organisation} Section~\ref{sec:related} positions the work. Section~\ref{sec:setting} gives the running example, the system and threat model and the problem statement. Sections~\ref{sec:model} to~\ref{sec:ranking} present the model, the admission of change, the protocol and the ranking. Section~\ref{sec:evaluation} reports the evaluation, Section~\ref{sec:discussion} discusses deployment and limitations, and Section~\ref{sec:conclusion} concludes.
